\cotefile{21}
\scanpage{2}
Notes pour Dieud.

\S~9
% archive sheet number: 1

\scanpage{4}
2. \uline{Généralités sur certains foncteurs de modules}.

Le présent numéro pourrait passer également au \deleted{\unsure{nº|Nº}} Chap $\unsure{0|O}_{\mathrm{III}}$, qui contiendrait aussi certains des sorites du par.7 .

\uline{Proposition} 2.1. Soient $\underline C$ une catégorie abélienne, $\underline G$ une partie de $\mathrm{Ob}\,\underline C$, telle que tout objet de $\underline C$ admette une suite de composition finie dont les \deleted{quotients}\inserted{facteurs} sont isomorphes à des quotients d'objets $M\handwritten{\in}\underline G$. \deleted{\unsure{Al}} Soit $\underline C'$ une seconde catégorie abélienne.

(i) Soit $T:\underline C\handwritten{\to}\underline{C'}$ un foncteur \deleted{add} exact à droite, pour que $T$ soit nul, il f et s qu'il soit nul sur les $M\handwritten{\in}\underline G$.

(ii) Soit $(T^i)$ un foncteur cohomologique de $\underline C$ dans $\underline C'$, tel que $T^i=0$ pour $i\handwritten{>}i_{\unsure{o|0}}$, et soit $n$ un entier. Pour qu'on ait $T^i=0$ pour $i\handwritten{>}n$, il f. et s. que pour tout $M\handwritten{\in}\underline G$, on ait $T^i(M)=0$ pour $i\handwritten{>}n$.

Dém. (i) est trivial, pour (ii) on procède par récurrence \inserted{descendante} sur $n$, c'est trivial pour $n\handwritten{\geq}i_o$, et pour $n$ quelconque, supposant que c'est fait pour $n+1$, on note que si les $T^i(M)$ ($M\handwritten{\in}\underline G$, $i\handwritten{>1+}n$) sont nuls, alors par l'hy de réc on a $T^i=0$ pour $i\handwritten{>}n+1$, donc $T^{n\handwritten{+1}}$ est exact à droite, et d'après (i) il est nul, cqfd.

\uline{Corollaire} \deleted{\unsure{3}}2.2. Sous les conditions de 2.1. soit $(T^i)$ un foncteur cohomologique contravariant de $\underline C$ dans $\underline C'$, tel que $T^i=0$ pour $i\handwritten{<}i_o$, et soit $n$ un entier. Pour qu'on ait $T^i=0$ pour $i\handwritten{<}\deleted{\unsure{m}}\handwritten{n}$, il f et s que pour tout $M\handwritten{\in}\underline G$, on ait $T^i(M)=0$ pour $i\handwritten{<}n$.

Il suffit d'appliquer 2.1. (ii) à $\underline C$ et ${\underline C'}^{o}$. Pour qu'il n'y ait pas de jaloux, il faudrait aussi répéter (i). Dans le présent paragraphe, ce sera sous cette forme que no\deleted{\unsure{j}}us utiliserons 2.1. (mais il aurait été canularesque d'énoncer 2.1. avec des foncteurs contravariants).

\uline{Corollaire} 2.3. Soient $\underline C$, $\underline C'$ des catégories a\deleted{\unsure{v}}béliennes, $\underline G$ une partie de $\mathrm{Ob}\,\underline C$, on suppose que les objets de $\underline C$ sont noethériens (i.e. satisfont la condition des chaines\handwritten{\unsure{/}}ascendantes pour les sous-objets), et que $\underline G$ est un système de générateurs de $\underline C$ (i.e. que pour tout $\deleted{\unsure{A|E}}F\handwritten{\ne}0$ dans $\underline C$, il existe un $M\handwritten{\in}\underline C$ et un homomorphisme non nul $M\handwritten{\to}F$). Alors les conclusions de 2.1. (i)(ii) et 2.2. (i)(ii) sont valables.

Il suffit de vérifier l\deleted{es}\handwritten{'} hypothèses de 2.1.\inserted{\handwritten{ sur $\underline G$}}, pour ceci soit $F\handwritten{\in}\mathrm{Ob}\,\underline C$, \deleted{\unsure{si F}} soit $F'$ un sous-objet maximal de $F$ admettant une suite de composition finie comme dite dans 2.1. (un tel $F'$ existe, car $0\handwritten{\subset}F$ satisfait à l'hypothèse de 2.1., et on applique l'hyp noeth sur $F$). Je dis que $F'=F$, en \deleted{\unsure{effe}}
% archive sheet number: 3

\scanpage{5}
effet sinon il existerait un homomorphisme non nul d'un $M\handwritten{\in}\underline G$ dans $F/F'$,
et l'image inverse \inserted{dans $F$} de son image \deleted{dans} contredirait l'hypothèse maximale
pour $F'$.

\uline{Exemple} 2.4. Soient $A$ un anneau comm noeth, $I$ un idéal de $A$, $X$ le spectr\unsure{e}
de $A$ et $Y$ la partie fermée de $X$ définie par $I$. Soit $\underline C$ la catégorie des
$A$-modules \inserted{de type fini} dont tout élément est annulé par une puissance de $I$, ou la catégorie équivalente des Modules \deleted{quasi-}cohérents sur $X$ dont le support est
contenu dans $Y$. Les objets de $\underline C$ sont noethériens. Si $\mathbb M\handwritten{\in}\mathrm{Ob}\,\underline C$, alors on
voit tout de suite que $\mathbb M$ \uline{est un générateur de} $\underline C$ \uline{sss} $\operatorname{supp}\mathbb M=Y$. C'est
nécessaire, car si on avait $\operatorname{supp}\mathbb M\handwritten{\ne}Y$, il existerait une composante
irréductible $Y_i$ de $Y$ non contenue dans\deleted{\unsure{s}} $\operatorname{supp}M$, et prenant $F=\underline O_{Y_i}$, avec
$Y_i$ muni de la structure réduite induite, on trouve que $\operatorname{Hom}(M,F)=0$, $F\handwritten{\ne}0$.
C'est suffisant, car si \deleted{supp} $F$ est \deleted{\unsure{selement|seulement}} quasi-cohérent, de support
$\handwritten{\subset}Y$, on a $\operatorname{Ass}\operatorname{Hom}(M,F)=\operatorname{supp}M\handwritten{\cap}\operatorname{Ass}F=\operatorname{Ass}F$, donc $\operatorname{Hom}(M,F)=0$ implique
$\operatorname{Ass}F=\handwritten{\varnothing}$ d'où $F=0$. \deleted{\unsure{(}}(Donc $M$ est un générateur également dans la catégorie de tous les m\deleted{\unsure{d}}odules \uline{quasi}-cohérents $F$ de support $\handwritten{\subset}Y$). \deleted{Par suite,
si $T$ pour tester la nullité de} Le même argument montre que si on a une
famille $(M_i)_{i\handwritten{\in}I}$ d'objets de $\underline C$, alors c'est un système de générateurs de $\underline C$
sss la \deleted{\unsure{re}}réunion des $\operatorname{supp}M_i$ est $Y$ (commencer par ça). Pour une famille
satisfaisant cette condition, l'hypothèse de 2.1. est vé\deleted{\unsure{r}}rifée, et par
suite, pour vérifier la nullité de foncteurs comme dans 2.1. ou 2.2.,
il suffit de tester par les $M_i$, \deleted{ou de tester \unsure{i}} en particulier il suffit
de tester par \uline{un} $M$ de support $Y$.

\uline{Exemple} 2.5. On peut généraliser 2.4. en prenant un schéma \deleted{\inserted{quasi-compact}}\deleted{$X$ muni d'un}
noethérien $X$, muni d'un faisceau ample noté $\underline O_X(1)$, \deleted{Soit $(M_i)$} et d'une
partie fermée $Y$, $\underline C$ désignant encore la catégorie des faisceaux cohérents
sur $X$ à support $\handwritten{\subset}Y$. Les objets de $\underline C$ sont encore noethériens, et si
$(M_i)_{i\handwritten{\in}I}$ est une famille d'objets de $\underline C$, telle que \deleted{$\operatorname{supp}M_{\unsure{Y|i}}\unsure{\subset|=}$}
\[
\handwritten{\bigcup_i}\operatorname{supp}M_i=Y\quad,
\]
(il aurait fallu déployer la formule é\handwritten{g}alement dans 2.4.) , alors les
$M_i(\deleted{\unsure{+}}\unsure{-|\pm}n)$, avec $n\handwritten{>}N$ ($N$ entier donné), forment un système de générateurs
de $\underline C$ (et même de la catégorie des Modules qu coh de support $Y$). En effet
si $F$ est un Module qu coh de support $Y$, et $F\handwritten{\ne}0$, alors \deleted{\unsure{l'arg}} on a vu
dans 2.4. que le faisceau $\underline{\operatorname{Hom}}(M_{\deleted{\unsure{n}}i},F)$ est non nul pour $i$ convenable (se placer dans un ouvert affine où $F$ est $\handwritten{\ne}0$), donc il existe des sections non
nulles de $\underline{\operatorname{Hom}}(M_i,F)(n)$ pour $n$ grand, donc des\handwritten{\unsure{/}}homomorphismes non nuls de
$M_i(-n)$ dans $F$ pour $n$ grand. -- On peut donc utiliser les $M_i(-n)$, $n\handwritten{>}N$,
% archive sheet number: 4

\scanpage{6}
comme dans 2.4., pour tester \handwritten{l}a nullité de foncteurs comme dans2.1.
et dans 2.2. D'ailleurs\deleted{\unsure{r}}, 2.4. est évidemment u\handwritten{n} cas particulier de 2.5.
en prenant $\underline O_X$ comme faisceau ample sur le schéma affine $X$.

Soient encore $A$ un anneau noethérien, $\underline C$ la catégorie des $A$-modules\deleted{\unsure{,}}
de type fini, $\underline C'$ une catégorie abélienne quelconque, et
\[
T:\underline C\handwritten{\longrightarrow}\underline C'
\]
un foncteur additif. Pour tout $F\handwritten{\in}\mathrm{Ob}\,\underline C$, on peut alors considérer $T(F)$
comme un $A$-$\underline C'$-module, i.e. un objet de $\underline C'$ muni d'un\deleted{\unsure{e}} homomorphisme de
$A$ dans l'anneau de ses endomorphismes, De cette façon, le foncteur $F$ se
factorise à travers la catégorie des $\handwritten{A}$-$\underline C'$-modules, et on peut dire encore
que la donnée de $T$ équivaut à la donnée d'un \uline{foncteur} \inserted{$A$-}\uline{linéaire} de $\underline C$ dans
la catégorie \deleted{$\underline C'(A)$} $A(\underline C')$ des $A$-$\underline C'$-modules (plus précisément, \deleted{\unsure{en|ce}} on a une
équivalence naturelle entre ces deux catégories de foncteurs ...). C'est
le sorite déjà développé au par.7, mais qu'il faudrait faire de cette
façon plus générale, pour pouvoir par exemple renverser les flêches. Pour
l'instant, le fait qu'on s'est limité dans $\underline C$ aux modules de type fini
n'a pas joué, n'importe quelle sous-catégorie \inserted{pleine} de la catégorie des $A$-modules
aurait fait l'affaire (on se sert simplement du fait que les objets de $\underline C$
sont munis \deleted{``fonctoriellement} chacun d'une structure de $A$-$\underline C$-module, et que
les homomorphimses dans $\underline C$ sont des $A$-homomorphismes ...). Nous voulons
maintenant définir un homomorphisme fonctoriel
\[
T(A)\otimes_A M\handwritten{\longrightarrow}T(M)\tag*{\handwritten{$(*)$}}
\]
dont la \deleted{définition}\inserted{construction} utilisera le fait que \deleted{$A$ \unsure{$\in$} $\mathrm{Ob}\,\underline C$, et} $\underline C$ contient $A$ et que
les objets de $\underline C$ sont des modules de présentation finie sur $A$ (assuré par
$A$ noethérien, et les objets de $\underline C$ de type fini). Tout d'abord, si $\underline C'$ est
une catégorie abélienne, $A$ un anneau quelconque (pas néc comm.), \deleted{et} $\deleted{\unsure{L|M}}\handwritten{M}$ un
$A$-module à gauche, $P$ un $A$-$\underline C'$-module à droite, on définit
\[
P\otimes_A M\handwritten{\in\mathrm{Ob}}\,\underline C'\quad,
\]
sous réserve d'existence, comme l'\deleted{unique} objet de $\underline C'$ (déterminé à isomorphisme unique près) qui représente le foncteur \margin{\handwritten{déployer}}%
% slanted
$Q\handwritten{\leadsto}\operatorname{Hom}_A(M,\operatorname{Hom}(P,Q))$ de
$\underline C'$ dans $(\mathrm{Ens})$, -- où\handwritten{,} dans le double $\deleted{\unsure{H}}\operatorname{Hom}$, le $\operatorname{Hom}(P,Q)$ est considéré
comme un $A$-module à gauche, grâce aux opérations de $A$ sur $P$ à droite. Donc
par définition on doit avoir un homomorphisme fonctoriel en $Q$
\[
\operatorname{Hom}_{\underline C'}(P\otimes_A M,Q)\handwritten{\simeq}\operatorname{Hom}_A(M,\operatorname{Hom}(P,Q))
\]
% archive sheet number: 5

\scanpage{7}
D'après les sorites généraux sur les foncteurs représentables, $P\otimes_A M$
dépend fonctoriellement de $\handwritten{(}M\handwritten{,P)}$ \deleted{et de $P$,} \deleted{là où il} pour les \inserted{\handwritten{\unsure{couples|couple} d'}}arguments pour
lesquels il est défini. De plus, on voit tout de suite que si on a une
suite exacte
\[
M''\handwritten{\to}M'\handwritten{\unsure{-|\to}}M\handwritten{\unsure{\to|-}}0
\]
de $A$-modules à gauche, telle que $P\otimes_A M'$ et $P\otimes_A M''$ existent, alors $P\otimes_A M$
existe et \deleted{on a une}\inserted{\handwritten{la}} suite \deleted{exacte}
\[
P\otimes_A M''\handwritten{\to}P\otimes_A M'\handwritten{\unsure{-|\to}}P\otimes_A M\handwritten{\unsure{-|\to}}0
\]
\inserted{\handwritten{est exacte}} (en particulier, $P\otimes_A M$ est \uline{foncteur exact à droite en} $M$, sous réserve
d'existence).
% The block beginning "On a le fait analogue" is enclosed and carried down by an arrow to "Combinant".
\deleted{Cela prouver} Plus généralement, le foncteur
$P\otimes_A M$ en $M$, ou en $P$, commute aux limites inductives quelconques (pas nécessairement filt\unsure{r}antes), au sesn suivant: si les $P\otimes_A M_i$ existent, \deleted{et $\lim$}
alors $P\otimes_A\mathop{\underset{\handwritten{\longrightarrow}}{\lim}}\limits_{\handwritten{i}}M_i$ existe sss $\underset{\handwritten{\longrightarrow}}{\lim}P\otimes_A M_i$ existe, et les deux sont can isomorphes. Ceci s'applique en particulier aux limites inductives sur un ensemble discret, i.e. aux sommes directes.
\inserted{On a le fait analogue pour des suites exactes en $P$\inserted{\handwritten{, et plus gén. des $\varinjlim$ en $P$}}. D'autre
part, $P\otimes_A A_g$ existe toujours et\handwritten{\unsure{/}}on a un isomorphisme canonique
\[
P\otimes_A A_g\handwritten{\simeq}P\;.
\]
\deleted{(Pour}(Tous ces faits sont triviaux sur les définitions, en interprétant
en termes de foncteurs en $Q$).}
Combinant ces remarques, on voit
que $P\otimes_A M$ existe \deleted{toujours} dès que $M$ est de présentation finie (car il
existe pour $M=A_g^{(I)}$, $I$ ensemble fini, donc pour des conoyaux d'homomorphismes $A_g^{(I)}\handwritten{\longrightarrow}A_g^{(J)}$, $I$ et $J$ finis), et sans restriction sur $M$ si dans $\underline C'$
les sommes directes quelconques existent, (même \deleted{\unsure{ar}}argument, sans restriction
de finitude sur $I,J$). De plus, on a un pro\deleted{\unsure{t}}cédé de construction de $P\otimes_A M$
par ``genérateurs et relations'' en termes de ``générateurs et relations''
pour $M$, i.e. en termes \deleted{d'un} de la matrice définissant une présentation
\deleted{finie} de $M$ (cette matrice définit une matrice sur $I\mathbin{\mathrm{x}}J$\handwritten{,} définissant une
application $P^{(I)}\handwritten{\longrightarrow}P^{(J)}$, dont le conoyau est le produit tensoriel cherché).
Tout ceci devrait figurer dans un numéro à part\inserted{ au Chap $0$}, sans trainer l'accent '
sur la catégorie abélienne. Du point de vue notations, si $C'$ est \deleted{la duale}\inserted{l'opposée}
d'une catégorie abélienne $C_1$, \deleted{on note $\operatorname{Hom}$} les $A$-$C'$-modules à droite
sont les $A$-$C_1$-modules à gauche, et le produit tensoriel dans $C'$ se note
(en termes de $C_1$)
\[
\operatorname{Hom}_A(M,P)
\]
défini directement par
\[
\deleted{\operatorname{Hom}_{Q}(\operatorname{Hom}(M,P))}
\]
% archive sheet number: 6

\scanpage{8}
\[
\operatorname{Hom}(Q,\operatorname{Hom}_A(M,P))\handwritten{\simeq}\operatorname{Hom}_A(M,\operatorname{Hom}(Q,P))\quad,
\]
(isomorphisme fonctoriel en l'objet variable $Q$ de $C_1$), où ce fois-ci $M$ \uline{et}
$P$ sont des $A$-modules à gauche ($P$ à valeurs dans $C_1$). On utilisera sans
référence spéciale les sorites duaux de ceux donnés pour $\otimes$, qui sont une
simple traduction de ces derniers. L'intérêt du traitement abst\deleted{\unsure{a}}rait donné
ici est précisément, entre autres, d'unifier certains énoncés \inserted{\deleted{\unsure{des}}\handwritten{sur \unsure{les|des}} modules habituels} portant, \deleted{\unsure{l'u}}
l'un sur des produits tensoriels, l'autre sur des $\operatorname{Hom}$,
(il faut avoir vu bien entendu qu'on retroyve \deleted{c}\handwritten{l}es \deleted{derniers lorsque}\inserted{\handwritten{\unsure{objets} \unsure{traditionnels|habituel} si}}
$C'$ resp. $C_1$ est la catégorie $(\mathrm{Ab})$, i.e. comme cas particulier de produits
tensoriels avec des $A$-$(\mathrm{Ab})$-modules et des $A$-$(\mathrm{Ab})^{o}$-modules\handwritten{)}.

Ces rappels étant faits, par définition la donnée d'un homomorphisme
$(\handwritten{*})$ équivaut à la donnée d'un $A$-homomorphisme
\[
M\handwritten{\longrightarrow}\operatorname{Hom}(T(A),T(M))\quad,
\]
or un tel s'obtient de façon naturelle en composant
\[
M\handwritten{\longrightarrow}\operatorname{Hom}_A(A,M)\handwritten{\longrightarrow}\operatorname{Hom}(T(A),T(M))
\]
Ceci posé, on a:

\uline{Proposition} 2.6. L'homomorphisme $(\handwritten{*}):T(A)\otimes_A M\handwritten{\to}T(M)$ est fonctoriel.
Pour que ce soit un isomorphisme, il f et s que $T$ soit\handwritten{\unsure{/}}exact à droite.

\deleted{Suffisance}\inserted{Nécessité} triviale, pour la suffisance, on note que tout homomorphisme \deleted{de} $T'(M)\handwritten{\to}T(M)$ de foncteurs exacts à droite, qui induit un isomorphisme $T'(A)\handwritten{\to}T(A)$, est un isomorphisme \deleted{pour} (car c'est un isomorphisme
pour un argument $L$ libre de type fini, donc aussi pour un conoyau
$L_1\handwritten{\to}L_{\unsure{o|0}}$ de tels arguments, puisque $T,T'$ sont exacts à droite ...).

\uline{Corollaire} 2.7. Le foncteur $T\handwritten{\leadsto}T(A)$ est une équivalence de la catégorie
\deleted{des} $\underline{\operatorname{Hom}}(\underline C,\underline C')_{\deleted{g}d}$ des foncteurs exacts à \deleted{gauche}\inserted{droite} de $\underline C$ \handwritten{d}ans $\underline C'$, dans\handwritten{\unsure{/}}la
catégorie $A(\deleted{\unsure{\mathrm{Ab}}}\handwritten{\underline C'})$ des $A$-\inserted{\handwritten{$\underline C'$-}}modules quelconques, une équivalence\handwritten{\unsure{/}}quasi-inverse
est le foncteur $H\handwritten{\leadsto}T_H$, où $T_H(M)=H\otimes_A M$.

\uline{Remarque} 2.8. L'homomorphisme $(\handwritten{*})$ est défini également, sans hypothèse
noethé\unsure{ri|fi}enne sur $A$, lorsque $\underline C$ est la catégorie de tous les $A$-modules,
p\handwritten{\unsure{our}}vu que dans $\underline C'$ les sommes directes quelconques existent (donc que
le premier membre ait un sens). Alors 2.6. \deleted{peut} reste valable en y
précisant: $T$ exact à droite \uline{et} commute aux sommes directes ; et de même
pour 2.7.

Nous auronségalement à utiliser la situation duale d'un foncteur
% archive sheet number: 7

\scanpage{9}
contravariant $T$ de $\underline C$ dans $\underline C'$, i.e. d'un foncteur
\[
T:\underline C^{o}\handwritten{\longrightarrow}\underline C'\quad,
\]
(ce cas se ramenant au précédent en remplaçant $\underline C'$ par $ {\underline C'}^{o}$). On trouve
maintenant un homomorphisme canonique, fonctoriel en $M$:
\[
T(M)\handwritten{\longrightarrow}\operatorname{Hom}_A(M,T(A))\quad,
\]
et il faudra expliciter les énoncés duals de 2.6. et 2.7. , en particulier
l'homomorphisme précédent est un isomorphisme sss $T$ est exact à gauche
(i.e. transforme une suite exacte $M''\handwritten{\unsure{-|\to}}M'\handwritten{-}M\handwritten{-}0$ dans $\underline C$ en une suite
exacte $0\handwritten{-}T(M)\handwritten{-}T(M')\handwritten{-}T(M'')$ ): on ob\handwritten{t}ient une \deleted{\unsure{é|e}}équivalence entre la
catégorie des\handwritten{\unsure{/}}foncteurs contravariants exacts à g de $\underline C$ dans $\underline C'$ et la
catégorie $\handwritten{A(}\underline C'\handwritten{)}$. Ces résultats s'étendent \inserted{\handwritten{comme dans 2.4.}} au cas où $\underline C$ est la catégorie de
tous les $A$-modules, sans hypothèse noethérienne sur $A$, pourvu que dans $\underline C'$
les produits quelconques existent.

Nous allons généraliser les considérations précédentes à la catégorie
$\underline C$ de l'exemple \deleted{\unsure{3}}2.4. Soit pour tout $n\handwritten{\geq}0$, $\underline C^n$ la sous-catégorie pleine de
$\underline C$ formée des modules de type fini sur $\underline C$ annulés par $I^{n+1}$, catégorie
équivalente à la catégorie des modules de type fini sur $A_n=A/I^{n+1}$. Alors
$\underline C_n\handwritten{\longrightarrow}\underline C$ est un foncteur exact, et $\underline C$ est réunion filtrante de ses sous-catégories $\underline C_n$. Alors un foncteur \handwritten{additif}
\[
T:\underline C\handwritten{\longrightarrow}\underline C'
\]
est défini par la suite des foncteurs restriction
\[
T_n:\underline C_n\handwritten{\longrightarrow}\underline C'\quad,
\]
et en vertu de ce qui précède, pour chaque $T_n$ on a un homomorphisme
fonctoriel
\[
T_n(A_n)\otimes_{A_n}M\handwritten{\longrightarrow}T_n(M)=T(M)
\]
qu'on peut aussi écrire
\[
\deleted{\unsure{T_n(}}\qquad T(A_n)\otimes_A M\handwritten{\to}T_{\deleted{\unsure{n}}}(M)\quad\handwritten{(n\geq n_{\unsure{o|0}})}\;.
\]
Pour $M$ fixé dans un $\underline C_{n_o}$, on obtient ainsi un système projectif \deleted{de}
\deleted{\unsure{$T_n$}} $(T(A_n)\otimes_A M)_{n\handwritten{\geq}n_o}$ d'objets de $\underline C'$, et un homomorphisme de ce dernier dans
le système projectif constant défini par \deleted{$M$} l'objet $T(M)$. Sous réserve
de l'existence de $\underset{\handwritten{\longleftarrow}}{\lim}T(A_n)\deleted{\unsure{\otimes M}}$, on trouve donc un homomorphisme can
\[
\underset{\handwritten{\longleftarrow}}{\lim}T(A_n)\otimes_A M\handwritten{\longrightarrow}T(M)
\]
d'où par composition un homomorphisme canonique
\[
(\underset{\handwritten{\longleftarrow}}{\lim}T(A_n))\otimes_A M\handwritten{\longrightarrow}T(M)\quad\handwritten{,}
\]
% archive sheet number: 8

\scanpage{10}
et dans ces derniers homomorphismes, il n'est plus nécessaire de faire
mention de l'entier $n_{\unsure{o|0}}$ (i.e. ils ne dépendent pas du choix de l'entier
$n_o$ tel que $M\handwritten{\in}\mathrm{Ob}\,\underline C_{n_o}$\handwritten{)}. Pour notre objet, nous aurons plutot à considérer la situation duale d'un foncteur contravariant sur $\underline C$ :
\[
T:\underline C^o\handwritten{\longrightarrow}\underline C'\quad,
\]
donnant lieu à des homomorphismes
\[
T_{\deleted{n}}(M)\handwritten{\longrightarrow}\operatorname{Hom}_A(M,T_{\deleted{\unsure{n}}}(A_n))\qquad(n\handwritten{\geq}n_o)\quad,\tag*{(\%)}
\]
satisfaisant les propriétés de transitivité qu'il faut pour pouvoir passer
à la lim\deleted{o}ite (sous réserve d'existence des $\underset{\handwritten{\longrightarrow}}{\lim}$ en question)
\[
T(M)\handwritten{\unsure{-|\to}}\underset{\handwritten{\longrightarrow}}{\lim}_{\handwritten{n}}\mathbb H\mathrm{om}_A(M,T(A_n))\handwritten{\to}\operatorname{Hom}_A(M,\underset{\handwritten{\longrightarrow}}{\lim}_{\handwritten{n}}T(A_n))
\]
Pour que $T$ soit exact à gauche, il f et s que les $T_n$ le soient, ou encore
que les homomorphismes (\%) soient des isomorphismes. Si dans $\underline C'$ les
limites inductives \inserted{filtrantes} dénombrables existent et sont ``exactes'', à fortiori
si l'axiome (Ab 5) est vérifié, alors \deleted{l'exactitude à gauche de $T$ impliqque donc que l'homomorphisme} on voit fa\handwritten{cil}ement que le foncteur
$P\handwritten{\leadsto}\operatorname{Hom}_A(M,P)$ de \deleted{$\underline C'$ dans $\underline C'$} $A(\underline C')$ dans $\underline C'$ commute aux limites inductives \inserted{filtr} dénombrables (pour tout $M$ de présentation finie sur $A$), Donc \deleted{sous} si
\inserted{\handwritten{de plus}} $T$ est exact à gauche, l'homomorphisme composé
\[
T(M)\handwritten{\longrightarrow}\operatorname{Hom}_A(M,\underset{\handwritten{\longrightarrow}}{\lim}_{\handwritten{n}}T(A_n))
\]
est un isomorphisme, la réciproque étant triviale. Il y a lieu de
résumer ce résultat dans une

\uline{Proposition} 2.9.

généralisant 2.8. sous les conditions dites, et en déduire le

\uline{Corollaire} 2.\deleted{9}\handwritten{10}. $\underline C$, $\underline C'$ étant comme dans 2.8. , la catégorie des
foncteurs exacts à \deleted{droite} gauche $T$ de $\underline C$ dans $\underline C'$ est équivalente à la
catégorie des $A$-$\underline C'$-modules $H$ qui sont ``localement annulés par des
puissances de $I$'', i.e. qui sont limites inductives des sous-modules
$H_n=\operatorname{Hom}_A(A_n,H)$. On obtient deux foncteurs quasi-inverses par $T\handwritten{\leadsto}\underset{\handwritten{\longrightarrow}}{\lim}_{\handwritten{n}}T(A_n\unsure{)}$
et par $H\handwritten{\leadsto}T_H$, où $T_H(M)=\operatorname{Hom}_A(M,H)$ .

\uline{Remarque} 2.1\deleted{0}\handwritten{1}. Lorsqu'on ne suppose rien sur l'existence de limites
inductives dans $\underline C'$, on obtient encore une équivalence entre $\underline{\operatorname{Hom}}(\underline C^o,\underline C')_g$
et la catégorie des systèmes inductifs ``co-$\deleted{\unsure{J}}\handwritten{I}$-adiques'' $(H_n)$ de $A$-\inserted{$C'$-}modules
i.e. des systèmes inductifs de tels modules tels que \inserted{\handwritten{a)}} chaque $H_n$ soit
annulé par $I^{n+1}$ (donc en fait un $A_n$-$\underline C'$-module) et b) pour $m\handwritten{\geq}n$, l'homo%
% archive sheet number: 9
\scanpage{11}%
morphisme de transition $H_n\handwritten{\longrightarrow}H_m$ induit un isomorphisme de $H_n$ sur le
sous-$A$-$\underline C'$-module $\operatorname{Hom}_A(A_n,H_m)$ des ``éléments annulés par $I^{n+1}$'' . Bien
entendu, sous l'hypothèse sur $\underline C'$ faite dans 2.\deleted{\unsure{8}}\handwritten{9}. et 2.10. on voit aussitôt directement que la catégorie de systèmes inductifs qu'on vient de
décrire, est équivalente à \deleted{celle menti} la catégorie de $A$-$\underline C'$-modules mentionnée dans 2.10.\handwritten{, qui est \unsure{donc} un cas \unsure{particulier} de 2.11.}

\uline{Corollaire} 2.12. Soient $A$ un anneau comm noeth. muni d'un idéal $I$,
$\underline C$ la catégorie des $A$-modules de type fini de support $\handwritten{\subset}V(I)$ i.e. annulés
par une puissance de $I$, $\underline C'$ une catégorie abélienne à limites inductives
filtrantes dénombrables exactes, $(T^i)$ un foncteur cohomologique
contravariant de $\underline C$ dans $\underline C'$, tel que $T^i=0$ pour $i\handwritten{<}i_{\unsure{o|0}}$, $n$ un entier,
$(M_j)$ une famille de $A$-modules de type fini tel que laréunion des
supports des $M_j$ soit égale à $Y=V(I)$. \deleted{Cond} Posons pour tout entier $i$:
\[
H^i=\underset{\handwritten{\longrightarrow}}{\lim}_{\!m}T^i(A_{\deleted{\unsure{n}}\handwritten{m}})\qquad(\text{ où }A_m=A/I^{m+1})\;.
\]
Alors les conditions suivantes sont équivalentes:

(i) On a $T^i=0$ pour $i\handwritten{<}n$.

(ii) On a $H^i=0$ pour $i\handwritten{<}n$

(iii) On a $T^i(M_j)=0$ pour $i\handwritten{<}n$ et tout $j$.

Dém. L'équivalence de (i) et (iii) est mise pour mémoire, et a été
vue dans 2.4. On a (i)$\handwritten{\Rightarrow}$(ii), pour l'implication inverse on
raisonne par récurrence sur $n$, c'est trivial pour $n\handwritten{\leq}i_o$, pour $n$ quelconque
et supposant que c'est fait pour les $n'\handwritten{<}n$, il reste à prouver que $T^{n\handwritten{-1}}=0$,
or $T^{n\handwritten{-1}}$ est exact à gauche, donc en vertu de 2.9. est isomorphe à
$\operatorname{Hom}_A(M,H^{n\handwritten{-1}})$, \inserted{\handwritten{donc nul,}} ce qui achève la démonstration. De plus, on voit:

\uline{Corollaire} 2.13. Si les conditions équivalentes (i) à (iii) de 2.12. sont
vérifiées, on a un isomorphisme fonctoriel
\[
T^{n\deleted{\unsure{-1}}}(M)=\operatorname{Hom}_A(M,H^{n\deleted{\unsure{-1}}})
\]

Il serait préférable d'inclure 2.13. dans l'énoncé de 2.12. et de
démontrer par récurrence (ii)$\handwritten{\Rightarrow}$(i) \uline{plus} l'énoncé final 2.13. , pour évi\deleted{\unsure{t}}ter des red\deleted{\unsure{q}}ites.
% archive sheet number: 10

\scanpage{12}
\uline{Corollaire} 2.14. $\underline C$ étant comme ci-dessus la catégorie des $A$-modules de type fini annulés par une puissance de $I$, soit $T:\underline C^o\handwritten{\to}(\mathrm{Ab})$ un foncteur
contravariant exact à gauche de $\underline C$ dans $(\mathrm{Ab})$, Soit $(Y_i)_{i\handwritten{\in}I}$ la famille des
composantes irréductibles de $Y\inserted{=V(I)}$, et soit $(Y_i)_{i\handwritten{\in}J}$ une sous-famille telle
que pour tout $M\handwritten{\in}\mathrm{Ob}\,\underline C$, la nullité de $T(M)$ (qui de toute façon ne dépend
que de $\operatorname{supp}M$ en vertu de 2.4.) soit équivalente à la relation ``$Y_i\handwritten{\not\subset}
\operatorname{supp}M$ pour tout $i\handwritten{\in}J$''. Soit $S\inserted{\handwritten{=(y_i)_{i\in J}}}$ l'ensemble des points génériques des
$Y_i$, $i\handwritten{\in}J$. Alors pour tout $M\handwritten{\in}\mathrm{Ob}\,\underline C$, on a
\[
\operatorname{Ass}T(M)=\operatorname{supp}M\handwritten{\cap} S\quad,
\]
en particulier le $A$-module $T(M)$ est sans idéaux premiers associés imm.

En effet, posant $H=\lim T(A/I^n)$, on a un isomorphisme
\[
T(M)\handwritten{\simeq}\operatorname{Hom}(M,H)
\]
d'où
\[
\operatorname{Ass}\deleted{\unsure{E|P}}T(M)=\operatorname{supp}M\handwritten{\cap}\operatorname{Ass}H\handwritten{,}\tag*{\%}
\]
et il reste à prouver la formule
\[
\operatorname{Ass}H=S\quad.
\]
\deleted{Soit donc} Evidemment $\operatorname{Ass}H\handwritten{\subset}\operatorname{supp}H\handwritten{\subset}Y$, il reste à prouver que si $y\handwritten{\in}Y$,
alors $y\handwritten{\in}\operatorname{Ass}H$ équivaut à $y\handwritten{\in}S$. Or soit $M=A/\underline p$, où $\underline p=\underline p_y$, \inserted{\handwritten{donc $V(M)=V(\underline p)=\bar y$,}} en
vertu de \% on a $T(M)\handwritten{\ne}0$ (i.e. $\operatorname{Ass}T(M)\handwritten{\ne\varnothing}$\handwritten{)} sss $\deleted{\unsure{V(\underline p)}}\handwritten{\bar y}$ \deleted{\unsure{$=\bar y$}} renoontre \deleted{pas} $\operatorname{Ass}H$,
or par hypothèse on a $T(M)\handwritten{\ne}0$ sss \deleted{$\operatorname{supp}M=V(\underline p)$}\inserted{\handwritten{$\bar y$}} \deleted{\unsure{ne}} rencontre \deleted{pas} $S$,
ou ce qui revient au même, sss $\deleted{\unsure{\underline p}}\handwritten{y\in}S$. Cela prouve d'abord que $\operatorname{Ass}H\handwritten{\subset}S$,
(prendre $y\handwritten{\in}\operatorname{Ass}H$), puis $S\handwritten{\subset}\operatorname{Ass}H$ puisque tout $y\handwritten{\in}S$ doit avoir une
spécialisation dans $S$, donc être dans $S$ puisque deux éléments de $S$ dont
\deleted{sans} l'un spécialise l'autre sont identiques.

Le résultat 2.10. nous incite à traduire les propriétés des
modules \inserted{$H$} de support $Y=V(I)$ en termes du foncteur exact à gauche
correspondant $M\handwritten{\leadsto}\operatorname{Hom}(M,\mathbb H)$ sur $\underline C$. Ici, nous nous contentons d'expliciter
le résultat suivant:

\uline{Proposition} 2.15. Avec les notations précédentes, pour que $H$ soit un
$A$-module injectif, il f et s que le foncteur $T_H(M)=\deleted{\unsure{H}}\operatorname{Hom}(M,H)$ sur la
catégorie $\underline C$ des $A$-modules de type fini \deleted{à supp\unsure{o}} annulés par une puissance
de $I$, soit exact à droite donc exact. En d'autres termes, 2.10. induit
une équivalence entre la catégorie des foncteurs contravariants \uline{exacts}
$T$ de $C$ dans$(\mathrm{Ab})$, avec la catégorie des $A$-modules injectifs à support $\handwritten{\subset}V(I)$
i.e. dont tout élément est annulé par une puissance de $I$.
% Faint inverted offset from the neighbouring typescript is visible beneath the final paragraph.
% archive sheet number: 11

\scanpage{13}
% The three lines on this face are rotated 180 degrees in the scan.
% The faint mirrored text is the writing of scan 12 showing through the sheet.
\uline{Corollaire} 2.14. Soient $A$ un anneau comm noeth, $I$ un idéal de $A$, $M$ un
$A$-module, $H_I^{\unsure{o|0}}(M)$ le sous-module de $M$ formé des éléments annulés par une
puissance convenable de $I$. Si $M$ est injectif
% archive sheet number: 12

\scanpage{14}
Nécessité triviale, prouvons la suffisance, i.e. si $M\handwritten{\leadsto}\operatorname{Hom}(M,\mathbb H)$ est
e\deleted{\unsure{d}}xact sur la sous-caté\handwritten{g}orie $\underline C$ de la catégorie des $A$-modules, alors ce
foncteur est exact partout. On sait (Tohoku) qu'il suffit de \deleted{le} vérifier
que pour tout sous-module $N$ de $M=A$, $\operatorname{Hom}(M,H)\handwritten{\to}\operatorname{Hom}(N,H)$ est surjectif,
donc on est ramené au cas où $M,N$ sont \deleted{surjectifs}\inserted{\handwritten{de type fini}}. Or un homom
\[
u:\mathbb N\handwritten{\to}H
\]
a comme image un sous-module de type fini de $H$, donc annulé par une
puissance de $I$, i.e. \inserted{\handwritten{$u$ est}} nul sur un sous-module $I^n\mathbb N$ de $\mathbb N$. En vertu de \uline{Krull},
$I^n\mathbb N$ contient une intersection $N\handwritten{\cap}I^mN$. Or on a une inclusion
\[
0\handwritten{\to}N/N\handwritten{\cap}I^mM\handwritten{\longrightarrow}M/I^mM
\]
\uline{dans} $\underline C$ , et $u$ se factorise par un homomorphisme
\[
u':N/N\; I^mM\handwritten{\longrightarrow}H\quad;
\]
ce dernier par hypothèse peut se prolonger à $M/I^mM$, d'où un homomorphisme
$v:M\handwritten{\to}H$ qui prolonge $u$, cqfd.

\uline{Corollaire} 2.16. Soit $\deleted{\unsure{M}}K$ un $A$-module, et soit $H_I^{\unsure{o|0}}(\deleted{M}K)$ le sous-module de $\deleted{M}K$
formé des éléments annulés par une puissance de $I$. Si $K$ est injectif, il
en est de même de $H_I^o(K)$.

En effet, lorsque $M$ parcourt $\underline C$, on a un isomorphisme fonctoriel
\[
\operatorname{Hom}(M,H_I^o(K))\handwritten{\overset{\sim}{\longrightarrow}}\operatorname{Hom}(M,K)\quad,
\]
de sorte qu'il suffit d'appliquer le critère 2.15.

\deleted{Si on rejette tout ce numéro (comme il me semble raisonnable)}
\uline{Remarques} 2.17. On peut généraliser 2.6. et 2.9. , ainsi que leurs
corollaires, au cas où on est sur un préschéma noethérien $X$, $\underline C$ désignant
la catégorie des Modules cohérents sur $X$ à support $\handwritten{\subset}Y$ ($Y=X$ dans le
cas \deleted{\inserted{dual}} de 2.6.), dans le cas d'un foncteur contravariant \inserted{exact à gauche} $T$ de $\underline C$ dans $(\mathrm{Ab})$:
un tel foncteur \deleted{se représente de façon} se représente (de façon essentiellement unique) comme $\handwritten{F\leadsto}\operatorname{Hom}(F,H)$, où $H$ est un Module quasi-cohérent sur
$X$ à support $\handwritten{\subset}Y$, \margin{\handwritten{déployer}}%
% slanted
encore donné par $\handwritten{H\simeq}\underset{\handwritten{\longrightarrow}}{\lim}_{\handwritten{n}}T(\underline O_X/\underline I^{\handwritten{\unsure{n}}})$. Pour le voir, on note
que $T$ est ``ind-représentable'' par le critère général donné dans A.Grothe\unsure{n}dieck, Technique de descente et théorèmes d'existence en Géom Alg II,
Sém Bourbaki ...., \handwritten{(}qui s'applique ici grâce au fait que les objets de
$\underline C$ sont noethériens\handwritten{),} et on note que $\handwritten{I}\mathrm{nd}\,(\underline C)$ est équivalente de façon naturelle à la catégorie de tous les Modules quasi-cohérents sur $X$ ... Nous
n'aurons sans doute pas besoin de ce fait dans ce Chapitre, où les résultats affines développés ici suffiront pour notre propos.
% archive sheet number: 13

\scanpage{15}
\deleted{\unsure{Fourbis|Fourbi} \unsure{généraux}
sur les \unsure{espaces} \unsure{topologiques},
en \unsure{particulier} \unsure{noethériens}}

profondeur $\geq 2$
pour \unsure{faisceaux|faisceau}
d'ensembles
% archive sheet number: 14

\scanpage{16}
Propriétés $(\unsure{P|\mathscr P}_1)$ et $(P_2)$ \unsure{pour} des faisceaux d'ensembles.

\deleted{\unsure{Proposition}}

\uline{\unsure{Définition} 1} \unsure{Soient|Soit} $X$ espace topologique, $Y$ une partie fermée
de $X$, \illegible{2} faisceau d'ensembles sur $X$. On dit
que $F$ \unsure{a} la propriété $(P_1)$ (resp. $(P_2)$)
\deleted{\unsure{rel. à}}\inserted{rel. à} \deleted{\unsure{un}} $Y$, si l'\unsure{hom} \unsure{canonique}
\[
F\longrightarrow i_*i^*(F)=i_*(F\mid U)
\]
(où $i:U=X-Y\longrightarrow X$ \unsure{est} l'\unsure{imm.} canonique)
est injective (resp. bijective).
\margin{N.B. \unsure{regarder|signaler} aussi la propriété: \unsure{surjective|surjection}....}
% slanted

(Cela signifie \unsure{donc} aussi que
pour tout ouvert $V\subset X$, l'application \deleted{\unsure{de}} restriction
\[
\Gamma(V,F)\unsure{-|\longrightarrow}\Gamma(V\cap U,F)
\]
est injective (resp. bijective).

\unsure{Appelons} \uline{espace noethérien local} un espace
noethérien dont \unsure{tte} partie fermée \unsure{irred} a un
pt générique et un seul\inserted{ [``espace noethérien \unsure{spectral|spécial}'']}, et tel que $X$ ait un
seul pt fermé. Définition de $\deleted{\unsure{X_x}}\operatorname{Loc}_x(X)$, $x$ pt d'un
espace loc. noethérien \unsure{spectral}. Définition de $\deleted{\unsure{F_x}}\operatorname{Loc}_x(F)$

\uline{Définition} 2 Soit $F$ un faisceau sur un espace
noethérien local $X$. On dit que $F$ \unsure{a} la
propriété $(P_1)$, resp. $(P_2)$, si \deleted{\unsure{le}} $F$ a la propriété
$(P_1)$ resp. $(P_2)$ rel. à $Y=\{x\}$, où $x$ est le pt fermé.
\margin{Cela signifie au \unsure{moins}, si $\deleted{\unsure{X}}U=X-\{x\}$, que $\Gamma(X,F)=F_x\longrightarrow\Gamma(U,\unsure{\Gamma|F})$ est injectif (resp. bijectif)}
% margin note written sideways
% Long vertical guides mark the two definitions. A bracket joins the sideways note to the foot of the page.
% archive sheet number: 15

